\PassOptionsToPackage{english,russian}{babel}
\documentclass[aspectratio=169]{beamer}
\usetheme{metropolis}

\usepackage[T2A]{fontenc}
\usepackage[utf8]{inputenc}
\usepackage{babel}  

\usepackage{listings}
\usepackage{xcolor}

\definecolor{codebg}{RGB}{248,248,248}
\definecolor{codeframe}{RGB}{220,220,220}
\definecolor{codekw}{RGB}{0,92,197}
\definecolor{codestr}{RGB}{163,21,21}
\definecolor{codecom}{RGB}{0,128,0}

\lstset{
basicstyle=\ttfamily\small,
backgroundcolor=\color{codebg},
frame=single,
rulecolor=\color{codeframe},
framesep=4pt,
keywordstyle=\color{codekw}\bfseries,
stringstyle=\color{codestr},
commentstyle=\color{codecom},
showstringspaces=false,
breaklines=true,
tabsize=2
}

\begin{document}



%------------------------------------------------

\begin{frame}{Технология распределенных реестров}

Технология распределенных реестров – это подход к обмену и хранению информации, при котором:

\begin{itemize}
\item каждый участник может обладать полноценной копией реестра
\item синхронизация копий реестра происходит на основе протокола достижения распределенного консенсуса, то есть соглашения среди участников на добавление новой информации
\item каждый участник взаимодействия может иметь доступ к истории транзакций
\end{itemize}

Такой подход позволяет обеспечить единый источник достоверных данных без необходимости централизованного доверенного посредника.

\end{frame}

%------------------------------------------------

\begin{frame}{Блокчейн}

Одной из наиболее известных реализаций распределенных реестров является блокчейн.

В блокчейне информация группируется в цблоки, которые связаны между собой в цепочку. Каждый новый блок содержит ссылку на предыдущий.

Это обеспечивает:

\begin{itemize}
\item целостность данных
\item невозможность изменения истории
\item прозрачность всех транзакций
\end{itemize}

 Термин <<\textbf{блокчейн}>>  иногда заменяет название всей технологии, потому что такой способ записи используется чаще других.

\end{frame}

%------------------------------------------------
\begin{frame}{Структура блокчейна}

\begin{center}
\includegraphics[width=\textwidth]{blockchain_structure.png}
\end{center}
\end{frame}

%------------------------------------------------

\begin{frame}{Типы распределенных реестров}

По уровню доступа участников сети распределенные реестры можно разделить на несколько типов.

\textbf{Публичные сети}

Любой участник может подключиться к сети и участвовать в ее работе.

\textbf{Частные сети}

Доступ к сети ограничен, участники известны и контролируются.

\textbf{Гибридные сети}

Часть данных является публичной, а часть доступна только определенным участникам.

Для корпоративных систем чаще всего используются частные или гибридные сети.

\end{frame}

%------------------------------------------------

\begin{frame}{Основные свойства распределенных реестров}

Технология распределенных реестров обладает рядом ключевых характеристик:

\begin{itemize}
\item децентрализация — данные хранятся на множестве независимых узлов в P2P-сети
\item неизменяемость — записанные данные невозможно удалить или изменить
\item консенсус — все узлы сети используют единый алгоритм для синхронизации и подтверждения правильности обновлений реестра
\item прозрачность — участники имеют доступ к истории операций
\item криптографическая безопасность — используются цифровые подписи
\item автоматизация — возможна реализация смарт-контрактов
\end{itemize}

\end{frame}

%------------------------------------------------

\begin{frame}{Смарт-контракты}

Смарт-контракт — это программа, хранящаяся в распределенном реестре и автоматически исполняющая условия договора.

Исполнение смарт-контракта происходит при наступлении определенных условий.

Например, в системе эксплуатации дронов смарт-контракт может проверять:

\begin{itemize}
\item наличие действующей сертификации дрона
\item наличие страхового полиса
\item наличие разрешения на выполнение полета
\end{itemize}

Если все условия выполняются, система может автоматически разрешить выполнение миссии.

\end{frame}

%------------------------------------------------

\begin{frame}{Платформы распределенных реестров}

Существует множество платформ для реализации распределенных реестров.

Наиболее известные из них:

\begin{itemize}
\item Ethereum — публичная блокчейн-платформа со смарт-контрактами на языке Solidity
\item Tendermint / Cosmos — модульная платформа для создания собственных блокчейн-сетей
\item Hyperledger Fabric — корпоративная платформа распределенного реестра
\end{itemize}

Для корпоративных систем с известными участниками наиболее подходящей является платформа Hyperledger Fabric.

\end{frame}

%------------------------------------------------

\begin{frame}{Почему выбран Hyperledger Fabric}

Hyperledger Fabric изначально проектировался для использования в корпоративной среде.

Основные преимущества:

\begin{itemize}
\item идентификация участников сети
\item контроль доступа к данным
\item высокая пропускная способность
\item низкая задержка подтверждения транзакций
\item возможность разделения данных между организациями
\end{itemize}

Эти характеристики делают платформу подходящей для построения инфраструктуры взаимодействия участников системы эксплуатации дронов.

\end{frame}

%------------------------------------------------

\begin{frame}{Архитектура Hyperledger Fabric}

Основными компонентами сети Hyperledger Fabric являются:

\begin{itemize}
\item Peer nodes — узлы сети, хранящие копии реестра
\item Ordering Service — сервис упорядочивания транзакций
\item Chaincode — смарт-контракты
\item Certificate Authority — центр сертификации участников
\item Channels — изолированные каналы обмена данными
\end{itemize}

Каждая организация в сети может иметь собственные узлы и участвовать в консенсусе.

\end{frame}

%------------------------------------------------

\begin{frame}{Типы узлов в сети}

В сети Hyperledger Fabric используются несколько типов узлов.

\textbf{Endorsing Peer}

Узел, который симулирует исполнение транзакции и возвращает результаты выполнения клиентскому приложению вместе со своей подписью.

\textbf{Ordering Service}

Распределенный сервис на нескольких узлах, служит для формирования новых блоков распределенного реестра и создания очередности исполнения транзакций.

\textbf{Committing Peer}

Узел, на котором происходит проверка транзакций на соответствие политике одобрения, установленной для конкретного приложения, перед их записью в реестр.

\end{frame}

%------------------------------------------------

\begin{frame}{Центр сертификации}

Для идентификации участников используется центр сертификации (Certificate Authority).

Он выполняет следующие функции:

\begin{itemize}
\item выдача цифровых сертификатов участникам сети
\item подтверждение принадлежности участников к организациям
\item обеспечение аутентификации пользователей и узлов сети
\end{itemize}

Сертификаты основаны на стандарте X.509 и инфраструктуре открытых ключей (PKI).

\end{frame}

%------------------------------------------------

\begin{frame}{Модель выполнения транзакций}

В большинстве блокчейн-платформ используется модель:

Order → Execute

Hyperledger Fabric использует модель:

Execute → Order → Validate

\begin{enumerate}
\item выполнение транзакции и ее одобрение
\item упорядочивание транзакций
\item проверка транзакций и запись в реестр
\end{enumerate}

Такая архитектура позволяет повысить масштабируемость и производительность сети.

\end{frame}

%------------------------------------------------

\begin{frame}{Взаимодействие приложений с сетью}

Приложения взаимодействуют с сетью Hyperledger Fabric через SDK.

Основные этапы выполнения транзакции:

\begin{enumerate}
\item приложение формирует транзакцию
\item транзакция отправляется endorsing peer
\item выполняется смарт-контракт
\item ordering service формирует блок транзакций
\item committing peers проверяют и записывают блок в реестр
\end{enumerate}

\end{frame}

%------------------------------------------------
\begin{frame}{Взаимодействие приложений с сетью}

\begin{center}
\includegraphics[width=\textwidth]{fabric_transaction_flow.png}
\end{center}
\end{frame}
%------------------------------------------------
\begin{frame}{Каналы}

Канал представляет собой изолированную подсеть внутри блокчейн-сети.

Каналы позволяют:

\begin{itemize}
\item разделять данные между организациями
\item обеспечивать конфиденциальность транзакций
\item реализовывать независимые бизнес-процессы
\end{itemize}

Каждый канал имеет собственный реестр и набор смарт-контрактов.

\end{frame}

%------------------------------------------------

\begin{frame}{Архитектура сети}

\begin{center}
\includegraphics[width=0.95\textwidth]{fabric_network_architecture.png}
\end{center}

\end{frame}
%------------------------------------------------
\begin{frame}[fragile]{Смарт-контракт сертификации прошивок}

В системе реализован отдельный смарт-контракт \textbf{FirmwareContract},
который отвечает за управление сертифицированными прошивками дронов.

\begin{lstlisting}[language=Go]
type FirmwareContract struct {
	contractapi.Contract
}
\end{lstlisting}

Контракт поддерживает следующие группы операций:

\begin{itemize}\setlength\itemsep{0.5em}
\item \textbf{создание записи:} \texttt{CertifyFirmware};
\item \textbf{чтение данных:} \texttt{ReadFirmware}, \texttt{ListFirmwares};
\item \textbf{обновление:} \texttt{UpdateFirmware};
\item \textbf{отзыв сертификации:} \texttt{RevokeFirmwareCertification}.
\end{itemize}

\end{frame}

\begin{frame}[fragile]{Структура объекта Firmware}


Основная сущность контракта --- объект \textbf{Firmware}, содержащий:

\vspace{0.5em}

\begin{lstlisting}[language=Go]
type Firmware struct {
    ID                 string   `json:"id"`
    SecurityObjectives []string `json:"security_objectives"`
    SoftwareObjectives []string `json:"software_objectives"`
    CertifiedAt        string   `json:"certified_at"`
    CertifiedBy        string   `json:"certified_by"`
}
\end{lstlisting}


\end{frame}

\begin{frame}[fragile]{Пример реализации FirmwareContract}
\begin{lstlisting}[language=Go]
func (c *FirmwareContract) CertifyFirmware(...) error {
    if err := assertRole(...); err != nil {
        return err
    }

    certifierID, _ := cid.GetID(ctx.GetStub())

    firmware := Firmware{...}

    key, err :=  makeFirmwareKey(ctx, id)
    if err != nil {
        return err
    }

    return putState(ctx, key, firmware)
}
\end{lstlisting}

\end{frame}

\begin{frame}{Логика работы контракта FirmwareContract}

Функция \texttt{CertifyFirmware} реализует создание новой записи
о сертифицированной прошивке в распределённом реестре.

При выполнении транзакции контракт:

\begin{enumerate}\setlength\itemsep{0.5em}
\item проверяет роль вызывающей стороны;
\item получает идентификатор участника из контекста Fabric;
\item формирует объект \texttt{Firmware};
\item создает составной ключ для записи в world state;
\item сохраняет объект в распределённом реестре.
\end{enumerate}


\end{frame}

\begin{frame}[fragile]{DronePropertiesContract}

Смарт-контракт для управления паспортами дронов и страховками.

\begin{lstlisting}[language=Go]
type DronePropertiesContract struct {
	contractapi.Contract
}
\end{lstlisting}

\textbf{Типы данных:}
\begin{itemize}
\item TypeCertificate — сертификаты типа
\item DronePass — паспорта дронов
\item InsuranceRecord — страховки
\end{itemize}

\textbf{Операции:}
\begin{itemize}
\item Выпуск сертификатов типа
\item Создание/чтение/обновление паспортов
\item Управление страховками
\end{itemize}

\end{frame}

\begin{frame}[fragile]{Смарт-контракт управления заказами}

Смарт-контракт \textbf{OrderContract} управляет заказами на полеты дронов.

\begin{lstlisting}[language=Go]
type OrderContract struct {
	contractapi.Contract
}
\end{lstlisting}

\textbf{Типы данных:}
\begin{itemize}
\item Order — заказ с распределением платежей
\item MissionInsurance — страховка миссии
\item FlightPermission — разрешение ОрВД
\item RestrictionViolation — нарушения
\end{itemize}

\textbf{Жизненный цикл заказа:}
Создание $\rightarrow$ Назначение $\rightarrow$ Одобрение $\rightarrow$ 
Подтверждение $\rightarrow$ Полёт $\rightarrow$ Завершение $\rightarrow$ Оплата

\end{frame}
\end{document}

